
% =====================================================================
\section{The Master Theorem, scope, and conclusion}
\label{sec:master}
% =====================================================================

We now show that the entire ladder T0--T8 is forced by, and only by, the
Finiteness Premise---and that the premise is itself forced by the identifiability
of any one thing in a non-completable whole (\Cref{sec:ground}). Each rung holds
for finite contact graphs and fails to be forced when finiteness or the positive
floor is dropped. This is the sense in which the structures of the paper are
necessary substructures of a finite world within an inexhaustible whole---and the
sense in which structure itself is the consequence of finitude meeting
non-completability.

\subsection{Where finiteness is used}

\begin{proposition}[Dependence on the premise]
\label{prop:dependence}
Each rung uses the Finiteness Premise \textup{(}\Cref{ax:finite}\textup{)} as
follows:
\begin{enumerate}[label=\textup{(T\arabic*)},leftmargin=3em,start=0]
\item the positive floor $\floorc>0$ on edge weights, with connectedness, forces
$\wt(\cut(U))\ge\floorc$ \textup{(}\Cref{thm:floor}\textup{)};
\item finiteness of $\V$ and the absence of a distinguished vertex force
determination by complementation \textup{(}\Cref{thm:negation}\textup{)};
\item the floor forces $\Res\ge\floorc$, and finiteness attains the minimising
partition, giving a positive region-valued residual
\textup{(}\Cref{thm:residual-region}\textup{)};
\item finiteness makes the isomorphism class and its partition lattice
well-defined, so $\Res$ is an attained invariant
\textup{(}\Cref{thm:residual-invariant}\textup{)};
\item the floor and finiteness localise to subgraphs, giving a positive,
attained private invariant \textup{(}\Cref{thm:private}\textup{)};
\item finite degree with branching forces a gate to break finitely many ties
\textup{(}\Cref{thm:gate-necessary}\textup{)};
\item discreteness of steps makes the committed count a strictly monotone natural
number \textup{(}\Cref{thm:nonreturn}\textup{)};
\item the quotient of a finite graph is a finite contact graph, inheriting the
floor, so a single gate is forced on a coherent collective
\textup{(}\Cref{thm:single-gate}\textup{)};
\item finiteness of the specification against the infinitude of the forward
closure forces open-endedness, and the monotone count forces non-duplication
\textup{(}\Cref{thm:open,thm:nondup}\textup{)}.
\end{enumerate}
\end{proposition}

\begin{proof}
Each clause restates the cited theorem's use of \Cref{ax:finite}; the proofs are
those given in \Cref{sec:floor,sec:negation,sec:residual,sec:invariant,%
sec:gating,sec:monotone,sec:quotient,sec:construct}.
\end{proof}

\subsection{Failure without finiteness}

We show the ladder is not forced once the premise is dropped, by exhibiting, for
each load-bearing structure, a non-finite or zero-floor graph in which it fails.

\begin{proposition}[Necessity of the premise]
\label{prop:necessity}
Drop either clause of \Cref{ax:finite}:
\begin{enumerate}[label=\textup{(\alph*)},leftmargin=2.2em]
\item \textbf{Zero floor.} If edge weights may approach $0$
\textup{(}$\inf_e\wt(e)=0$\textup{)}, then a part may be separated from the rest
by a cut of arbitrarily small weight; $\wt(\cut(U))$ is no longer bounded below,
so T0 fails, and with it the residual positivity of T2 and the gate-tie
guarantees that rely on a floor.
\item \textbf{Infinite vertex set.} If $\abs{\V}=\infty$, a vertex may have
infinite degree, so a gate must break infinitely many ties at once and the finite
tie-breaking of T5 is not available in the same form; minimising partitions need
not be attained, so the residual of T2 need not be realised; and a walk's
``committed count'' may fail to be a well-ordered monotone natural number if
steps are not discrete, weakening T6.
\item \textbf{Complete graph with vanishing weights.} On an infinite complete
graph with weights tending to $0$, every part is adjacent to every other at
negligible cost: there is no positive residual, no privileged complement
structure, and no forced gate---none of T0--T8 is compelled.
\end{enumerate}
Hence the substructures are consequences of the premise, not of graphs in
general.
\end{proposition}

\begin{proof}
(a) Take a graph with a sequence of cuts of weight $2^{-n}\to0$ separating a
fixed part; then $\inf_U\wt(\cut(U))=0$, contradicting the conclusion of
\Cref{thm:floor}, which used $\wt(e)\ge\floorc>0$. (b) On an infinite star, the
centre has infinite degree, so no single sub-normalised distribution over its
incident edges selects after finitely many comparisons; and the partition lattice
of an infinite vertex set need not attain its infimum, so $\Res(\Gph)$ as a
minimum may not exist. (c) On $K_\infty$ with $\wt(e_n)=2^{-n}$, for any part $U$
one may route its separation through ever-cheaper edges so that
$\inf\wt(\cut(U))=0$; the floor, the positive residual, and the negation-based
individuation (which needs a determinate complement structure) all dissolve.
Each construction violates the conclusion of the corresponding rung, so finiteness
with a positive floor is necessary.
\end{proof}

\subsection{The Master Theorem}

\begin{theorem}[Structure is the shadow of finitude]
\label{thm:master}
Let $\Gph$ be a graph. The substructures T0--T8 are all forced if $\Gph$ is a
finite contact graph \textup{(}\Cref{ax:finite}\textup{)}, and none is forced if
$\Gph$ is infinite or admits edge weights with infimum $0$. Consequently a graph
is necessarily structured---carries the floor, the negation-determined parts, the
conserved region-valued residual, the private invariants, the gated trajectories,
the non-returning histories, the single-gated collectives, and the bounded
authorability---\emph{if and only if} it is finite with a positive floor. Structure
is neither imposed on a world from outside nor present in it independently of its
parts; it is exactly the form a world of finitely many positively separated parts
is compelled to take.
\end{theorem}

\begin{proof}
The forward direction is \Cref{prop:dependence} together with the theorems it
cites: under \Cref{ax:finite} each of T0--T8 holds. The reverse direction is
\Cref{prop:necessity}: dropping finiteness or the positive floor defeats each
load-bearing rung by explicit construction. Hence the conjunction of T0--T8 is
equivalent, over graphs, to the Finiteness Premise: finite-with-floor on one
side, the full ladder of structure on the other. The closing statement is the
biconditional read in words.
\end{proof}

\subsection{Identifiability and the trace of the whole}

The Master Theorem ties the ladder to the Finiteness Premise. We close the
deduction by tying that premise, in turn, to the bare fact of identifiability,
using \Cref{sec:ground}.

\begin{theorem}[Identifiability Theorem]
\label{thm:identifiability}
Let $\med$ be a whole. The following are equivalent:
\begin{enumerate}[label=\textup{(\roman*)},leftmargin=2.2em]
\item there exists an identifiable thing in $\med$---a proper part told from the
rest of a non-completable whole;
\item the contact graph of $\med$'s parts satisfies the Finiteness Premise with a
strictly positive floor \textup{(}\Cref{ax:finite}\textup{)};
\item the full ladder T0--T8 holds of that graph.
\end{enumerate}
Consequently a world in which T0--T8 fail is one in which nothing is
identifiable---only the undifferentiated whole. Structure and identifiability are
the same fact.
\end{theorem}

\begin{proof}
$(\mathrm{i})\Rightarrow(\mathrm{ii})$: an identifiable thing is a proper part
$A\neq\med$ compared against $\med\setminus A$; by \Cref{thm:floor-infinitude} the
comparison cannot complete and $\idres(A)>0$, and by \Cref{cor:floor-derived} the
infimum of these residuals is a positive floor $\floorc$, which is the edge-weight
floor of the contact graph of $\med$'s parts---the Finiteness Premise.
$(\mathrm{ii})\Rightarrow(\mathrm{iii})$: the Master Theorem (\Cref{thm:master}),
forward direction. $(\mathrm{iii})\Rightarrow(\mathrm{i})$: T0 alone provides a
part with a positive boundary cut (\Cref{thm:floor}), i.e. a proper part told from
the rest at positive cost---an identifiable thing. The contrapositive of the cycle
gives the closing statement: no identifiable thing $\iff$ no positive floor
$\iff$ the ladder fails, leaving only the whole undifferentiated.
\end{proof}

\begin{corollary}[The floor is the trace of the whole on its parts]
\label{cor:trace}
The positive floor is neither a measurement artifact nor a free constant: it is
the strictly positive identification residual that a non-completable whole forces
on each of its finite parts \textup{(}\Cref{thm:floor-infinitude}\textup{)}. Were
the whole completable, the residual could vanish, a thing's truth could be a
point, and the ladder would not be forced; because the whole is non-completable,
the residual is positive, truth is a cell, and the ladder stands. The floor is the
trace the inexhaustible whole leaves on whatever is finite within it.
\end{corollary}

\begin{proof}
\Cref{thm:floor-infinitude} and \Cref{cor:floor-derived} give the positive
residual; \Cref{prop:necessity} gives the failure of the ladder when the floor is
absent (completable or zero-floor case). The contrast is the stated trace.
\end{proof}

\subsection{Scope}

The theorems are statements about finite weighted graphs and an embedded gating
datum; nothing more is assumed. The single premise (\Cref{ax:finite}) is a finite
graph whose edges cost something. We have made no claim about any particular
system. Where each rung uses the premise is marked (\Cref{prop:dependence}); that
the premise is necessary is shown by counterexample (\Cref{prop:necessity}).
Interpretive readings---of parts as agents, gates as selecting fields, quotient
gates as collective purpose, residuals as conserved truth, the medium as the
whole, private invariants as identities, measurement as inquiry, authored
structures as manufactured individuals---are confined to remarks
(\Cref{rem:negation-interp,rem:repacking,rem:character,rem:field,%
rem:purpose,rem:onelaw,rem:new-category,rem:floor-trace,rem:marbles,%
rem:identifiability,rem:instrument-reading}) on which no theorem depends. A reader
will recognise that finitely many bounded parts---of many kinds---satisfy the
premise; we prove everything of the abstract graph alone, and material that could
not be proved at this standard has been omitted. One hypothesis beyond the graph
is used, and only in \Cref{sec:ground}: that the whole is non-completable. It is
not a cardinality posit but an order condition (no terminal stage), and it is
exactly what forces the floor; absent it, the floor may vanish and the ladder is
not compelled.

\subsection{Conclusion}

From the single hypothesis that a world is a finite graph of positively separated
parts, an entire architecture follows of necessity: no separation is free; parts
are fixed by what they are not; the boundary content between parts is a conserved
region that no relabelling removes and no point realises; each part carries a
private conserved invariant beneath its changing trajectories; those trajectories
require a selector and never return to a prior state; a coherent collective of
parts has exactly one selector on its quotient; and such structures, being
completely specifiable, can be authored---though authorship reaches only the
instrument and not its music, produces only new individuals and never copies, and
resolves only to its own floor. And beneath the hypothesis lies a barer one: that anything is identifiable at all
within a whole one cannot exhaust. From that alone the positive floor is
forced---the comparison identifying a thing against the rest can never finish, so
an irreducible residual remains---and with the floor, the whole ladder. To
measure a thing is then to reshuffle the knowledge around it, conserving the whole
and tightening the boundary, reducing how much is unsettled about the thing
without ever delivering the thing: identity is the conserved interior no
reshuffling reaches, and to resolve it to a point would be to exhaust the whole.
Each result is forced by finitude meeting the inexhaustible and dissolves without
either. The answer to why a world is structured at all is therefore not that
structure is added to it or discovered in it, but that structure is what a world
\emph{must} be for finitely many parts to stand apart within a whole they cannot
exhaust. The floor is the trace of the whole on its parts; structure is the shadow
of finitude meeting the inexhaustible; and identifiability, conserved truth, and
structure are one fact. The contact graph is its form.

% =====================================================================
\bibliographystyle{plainnat}
\bibliography{references}
% =====================================================================

\end{document}
